\documentclass[11pt,reqno]{amsart}
\usepackage[T1]{fontenc}
\usepackage[utf8]{inputenc}
\usepackage{lmodern}
\usepackage{amsmath,amssymb,amsfonts,amsthm,mathtools,mathrsfs,bm}
\usepackage{enumitem}
\usepackage{booktabs}
\usepackage{microtype}
\usepackage[colorlinks=true,linkcolor=blue,citecolor=blue,urlcolor=blue]{hyperref}
\usepackage[nameinlink,capitalize]{cleveref}
\allowdisplaybreaks
\setlength{\emergencystretch}{2em}
\newtheorem{theorem}{Theorem}[section]
\newtheorem{proposition}[theorem]{Proposition}
\newtheorem{lemma}[theorem]{Lemma}
\newtheorem{corollary}[theorem]{Corollary}
\theoremstyle{definition}
\newtheorem{definition}[theorem]{Definition}
\newtheorem{assumption}[theorem]{Assumption}
\newtheorem{example}[theorem]{Example}
\theoremstyle{remark}
\newtheorem{remark}[theorem]{Remark}
\newcommand{\E}{\mathbb E}
\newcommand{\Pp}{\mathbb P}
\newcommand{\Qp}{\mathbb Q}
\newcommand{\R}{\mathbb R}
\newcommand{\N}{\mathbb N}
\newcommand{\T}{\mathbb T}
\newcommand{\cE}{\mathcal E}
\newcommand{\cP}{\mathcal P}
\newcommand{\cL}{\mathcal L}
\newcommand{\Id}{\mathrm{Id}}
\newcommand{\tr}{\operatorname{tr}}
\newcommand{\Var}{\operatorname{Var}}
\newcommand{\Cov}{\operatorname{Cov}}
\newcommand{\osc}{\operatorname{osc}}
\newcommand{\TV}{\mathrm{TV}}
\newcommand{\norm}[1]{\left\lVert #1\right\rVert}
\newcommand{\abs}[1]{\left\lvert #1\right\rvert}
\newcommand{\pair}[2]{\left\langle #1,#2\right\rangle}

\title[Specular path ensembles and conjugate fields]{Geometry-Selected Path Ensembles and Endogenous Conjugate Fields for Specular Open Billiards}
\author{Qian Qi}
\address{Peking University, Beijing 100871, China}
\email{qiqian@pku.edu.cn}
\subjclass[2020]{Primary 37D50, 60F10; Secondary 37C30, 37A50}
\keywords{specular open billiards, path ensembles, pressure, large deviations, Gibbs conditioning, Doob transform}
\date{August 29, 2026}
\begin{document}
\raggedbottom
\begin{abstract}
\noindent\textbf{Platform: }\texttt{OB3-PE-v1}.\par
We construct a deterministic path-ensemble theory on one genuine planar
specular billiard.  Three strictly convex circular obstacles satisfy the
no-eclipse condition uniformly under a radial deformation.  The trapped
collision dynamics is conjugate to one mixing subshift of finite type, and the
actual free-flight time is a positive Holder roof.  The base path law is not an
arbitrary transition matrix: it is the Bowen--Margulis equilibrium state
selected by the billiard roof through the zero-pressure equation.

The mechanical path observable is the collision frequency at the moving
obstacle.  Its pressure is analytic and strictly convex because two physical
periodic billiard orbits have different collision frequencies.  An imposed
macroscopic impact frequency $r$ therefore selects a unique conjugate field
$\theta=I'(r)$.  A spectral local limit argument gives a quantitative
finite-window Gibbs-conditioning theorem.  The canonical driven preparation
is the equilibrium state of the tilted geometric potential on the same
specular billiard; no new stochastic transition law or surrogate map is
inserted.

The normalized Ruelle operator yields an exact Doob-normalized excess-pressure
semigroup.  This is the microscopic theta-evaluation used in the later papers.
The result upgrades the exact full-cover v6 mechanism to a real specular
billiard while retaining one typed deterministic platform.
\end{abstract}
\maketitle
\tableofcontents

\section{The explicit specular family}\label{sec:geometry}
For
\[
 \rho\in I_\rho:=\left[-\frac1{20},\frac1{20}\right]
\]
let
\[
 K_1(\rho)=\overline B((0,0),1+\rho),\qquad
 K_2=\overline B((6,0),1),\qquad
 K_3=\overline B((0,9),1).
\]
The billiard domain is the exterior of the three obstacles, the particle has
unit speed, and every collision obeys the Euclidean equal-angle specular law.

\begin{proposition}[Uniform no-eclipse geometry]\label{prop:no-eclipse}
For every $\rho\in I_\rho$ the obstacles are disjoint and satisfy the
no-eclipse condition.  The pairwise separation margins are at least
\[
 \frac{79}{20},\qquad \frac{139}{20},\qquad 3\sqrt{13}-2,
\]
and the three no-eclipse margins are at least
\[
 \frac{139}{20},\qquad \frac{79}{20},\qquad
 \frac{18}{\sqrt{13}}-\frac{41}{20}>0.
\]
\end{proposition}
\begin{proof}
The center distances are $6$, $9$, and $\sqrt{117}=3\sqrt{13}$, while the
largest first radius is $21/20$.  This gives the separation bounds.  Moreover
\[
 \operatorname{conv}(B(c_i,r_i),B(c_j,r_j))
 \subset[c_i,c_j]+B(0,\max\{r_i,r_j\}).
\]
The distances of the third center to the three center segments are $9$, $6$,
and $18/\sqrt{13}$.  Subtracting the relevant radii gives the displayed
strictly positive margins.
\end{proof}

Let $M_\rho$ be the non-wandering collision set and $B_\rho:M_\rho\to M_\rho$
the collision map.  Put
\[
 \Sigma=\{\omega\in\{1,2,3\}^{\mathbb Z}:\omega_{k+1}\ne\omega_k\},
 \qquad \sigma\omega=(\omega_{k+1})_{k\in\mathbb Z}.
\]

\begin{theorem}[Fixed symbolic desingularization]\label{thm:coding}
There are coding homeomorphisms $\chi_\rho:\Sigma\to M_\rho$ satisfying
\[
 B_\rho\circ\chi_\rho=\chi_\rho\circ\sigma.
\]
For every prescribed finite $N$, the collision points and the free-flight roof
\[
 \tau_\rho(\omega)
 =\abs{q_1(\rho,\omega)-q_0(\rho,\omega)}
\]
are $C^N$ in $\rho$ with values in a common Holder space, after reducing the
Holder exponent if necessary.  There are constants
$0<\tau_*\le\tau^*<\infty$ independent of $\rho$.
\end{theorem}
\begin{proof}
No eclipse gives a unique trapped orbit for each admissible obstacle code and
a uniform graph-transform contraction on exponentially weighted collision
sequences.  The margins in \cref{prop:no-eclipse} make all constants uniform.
Differentiation of the fixed-point equation gives the parameter regularity.
Uniform non-grazing and compactness give the roof bounds.  This is the standard
open-billiard construction; see \cite{ChernovMarkarian2006,Wright2016}.
\end{proof}

\begin{theorem}[Marked-period geometric nontriviality]\label{thm:marked-period}
The radial deformation is not a constant-time rescaling or a cohomologically
trivial roof perturbation.  The two-collision orbit alternating between
obstacles $1$ and $2$ has physical period
\[
 L_{12}(\rho)=2\bigl(6-(1+\rho)-1\bigr)=8-2\rho,
\]
whereas the orbit alternating between obstacles $2$ and $3$ has period
\[
 L_{23}=2(\sqrt{117}-2),
\]
independent of $\rho$.  Consequently $\partial_\rho\tau_\rho$ is not
cohomologous to a constant on the collision shift.
\end{theorem}
\begin{proof}
Both period-two trajectories lie on the line joining the corresponding disk
centres and obey the specular law normally.  Their lengths are the displayed
twice-gap formulas.  If $\partial_\rho\tau_\rho=c+U-U\circ\sigma$, its sum on
every period-two orbit would be $2c$.  Differentiating the two marked periods
gives respectively $-2$ and $0$, a contradiction.
\end{proof}

Choose a one-sided Holder representative $\widehat\tau_\rho$ cohomologous to
$\tau_\rho$.  The suspension defined by either roof is flow conjugate after a
bounded change of the time coordinate.

\section{The geometry-selected base ensemble}\label{sec:base-ensemble}
Let $P_\sigma(\varphi)$ denote topological pressure on the one-sided version
$\Sigma^+$.  Since $\widehat\tau_\rho>0$, there is a unique $h_\rho>0$ such
that
\begin{equation}\label{eq:bowen-root}
 P_\sigma(-h_\rho\widehat\tau_\rho)=0.
\end{equation}
Put
\[
 \varphi_\rho=-h_\rho\widehat\tau_\rho.
\]

\begin{theorem}[Bowen--Margulis preparation]\label{thm:BM}
The potential $\varphi_\rho$ has a unique equilibrium state $\nu_{\rho,0}$.
Its natural extension pushed forward by $\chi_\rho$ is an invariant
probability $\mu_{\rho,0}$ on the trapped specular billiard.  The suspension
measure is the unique measure of maximal entropy for the billiard flow.
All these objects are selected by the billiard geometry and the variational
principle; no transition matrix is prescribed.
\end{theorem}
\begin{proof}
The shift is topologically mixing and $\varphi_\rho$ is Holder, so the Ruelle
operator has a simple positive leading eigenvalue and a unique equilibrium
state \cite{Bowen1975,Ruelle1978,ParryPollicott1990}.  Equation
\eqref{eq:bowen-root}, Abramov's entropy formula, and the variational principle
identify the suspension state as the measure of maximal entropy.
\end{proof}

\begin{remark}[Preparation boundary]
The trapped set has zero ambient Liouville volume.  The theorem uses the
geometry-selected Bowen--Margulis preparation on the non-wandering set, not a
claim that ordinary incoming Lebesgue data are already distributed according
to $\mu_{\rho,0}$.
\end{remark}

\section{A mechanical impact current}\label{sec:current}
Define the collision-count observable
\[
 J(\omega)=\mathbf1_{\{\omega_0=1\}}.
\]
Thus $S_nJ$ is the actual number of impacts on the radially moving obstacle in
the first $n$ collisions.  Define
\begin{equation}\label{eq:Lambda}
 \Lambda_\rho(\theta)
 =P_\sigma(\varphi_\rho+\theta J)-P_\sigma(\varphi_\rho)
 =P_\sigma(\varphi_\rho+\theta J).
\end{equation}
Let $\nu_{\rho,\theta}$ be the equilibrium state of
$\varphi_\rho+\theta J$.

\begin{theorem}[Strict pressure duality]\label{thm:pressure-duality}
The function $\Lambda_\rho$ is real analytic and
\[
 \Lambda_\rho'(\theta)=\int J\,d\nu_{\rho,\theta},
\]
\[
 \Lambda_\rho''(\theta)
 =\lim_{n\to\infty}\frac1n
 \Var_{\nu_{\rho,\theta}}(S_nJ)>0.
\]
Moreover $\Lambda_\rho'$ is a real-analytic increasing bijection from
$\R$ onto $(0,1/2)$.  For every prepared impact frequency $r\in(0,1/2)$ there
is a unique
\[
 \theta_\rho(r)=(\Lambda_\rho')^{-1}(r).
\]
The large-deviation rate
\[
 I_\rho(r)=\sup_{\theta\in\R}
 \{\theta r-\Lambda_\rho(\theta)\}
\]
satisfies
\[
 I_\rho'(r)=\theta_\rho(r),\qquad
 I_\rho''(r)=\frac1{\Lambda_\rho''(\theta_\rho(r))}>0.
\]
\end{theorem}
\begin{proof}
Analytic perturbation of the Ruelle operator gives the derivative identities.
If the variance vanished, $J-\int Jd\nu_{\rho,\theta}$ would be a Holder
coboundary.  The period-two billiard orbit alternating between obstacles $2$
and $3$ has impact frequency zero at obstacle $1$, while the period-two orbit
alternating between $1$ and $2$ has frequency $1/2$.  Periodic-orbit sums of a
coboundary cannot have two different averages, proving strict positivity.
No admissible orbit visits obstacle $1$ consecutively, hence every invariant
measure has frequency at most $1/2$; the two displayed periodic measures attain
the endpoints $0$ and $1/2$.  Convex pressure asymptotics give the range of
$\Lambda_\rho'$.  Legendre duality gives the final formulas
\cite{Kifer1990,DemboZeitouni1998}.
\end{proof}

\section{Spectral Gibbs conditioning}\label{sec:conditioning}
For $n\ge1$ put
\[
 K_n^{\leftrightarrow}=\sum_{k=-n}^{n-1}J\circ\sigma^k.
\]
An integer sequence $k_n$ is called admissible if
$\{K_n^{\leftrightarrow}=k_n\}$ has positive $\nu_{\rho,0}$-probability.
The symmetric window is essential: a cylinder at one endpoint of a conditioned
block carries a right- or left-eigenfunction boundary bias, whereas a fixed
window in the bulk converges to the stationary tilted equilibrium state.

\begin{theorem}[Central-window microcanonical equivalence]\label{thm:gibbs-conditioning}
Fix compact sets $K_\rho\Subset I_\rho$ and
$K_r\Subset(0,1/2)$.  Let $k_n/(2n)\to r\in K_r$ and put
$\theta=\theta_\rho(r)$.  For every cylinder $C$ depending only on coordinates
$-\ell,\ldots,\ell$, there is $C_{K_\rho,K_r,\ell}$ such that, for admissible
$k_n$ with $\abs{k_n-2nr}=O(1)$,
\[
 \abs{
 \nu_{\rho,0}(C\mid K_n^{\leftrightarrow}=k_n)
 -\nu_{\rho,\theta}(C)}
 \le \frac{C_{K_\rho,K_r,\ell}}{\sqrt{n-\ell}}.
\]
Consequently every fixed central collision window under the microcanonical
impact preparation converges in total variation to the stationary canonical
equilibrium state on the same specular billiard.
\end{theorem}
\begin{proof}
Consider the Fourier family
\[
 \cL_{\rho,\theta+it}f(x)
 =\sum_{\sigma y=x}
 e^{\varphi_\rho(y)+(\theta+it)J(y)}f(y).
\]
Near $t=0$ analytic perturbation gives a simple eigenvalue with logarithm
\[
 \Lambda_\rho(\theta+it)
 =\Lambda_\rho(\theta)+itr
 -\frac12\Lambda_\rho''(\theta)t^2+O(t^3).
\]
Away from $2\pi\mathbb Z$ the normalized spectral radius is strictly below
one.  Indeed the period-two cycles $23$ and $12$ have $J$-sums $0$ and $1$;
therefore an equality-case multiplicative coboundary forces $e^{it}=1$.
Fourier inversion gives the lattice local limit theorem.  To evaluate the
central cylinder, propagate $n-\ell$ steps from both sides of $C$.  The leading
left and right eigenvectors then occur in matched pairs and produce the
stationary density $h_{\rho,\theta}\ell_{\rho,\theta}$; all boundary-biased
terms are separated from the central window by $n-\ell$ iterates and decay
spectrally.  Dividing by the denominator local-limit expansion gives the
stated ratio and rate, uniformly on the compact parameter sets.  This is the
Nagaev--Guivarc'h method with a bulk insertion \cite{HervePene2010}.
\end{proof}

\section{The same-billiard driven preparation}\label{sec:doob}
Let $e^{\Lambda_\rho(\theta)}$, $h_{\rho,\theta}>0$, and
$\ell_{\rho,\theta}$ be the leading eigenvalue, right eigenfunction, and left
eigenfunctional of $\cL_{\rho,\theta}$, normalized by
$\ell_{\rho,\theta}(h_{\rho,\theta})=1$.  Define
\begin{equation}\label{eq:markov-op}
 \cP_{\rho,\theta}f
 =e^{-\Lambda_\rho(\theta)}
 h_{\rho,\theta}^{-1}
 \cL_{\rho,\theta}(h_{\rho,\theta}f).
\end{equation}

\begin{proposition}[Doob-normalized equilibrium dynamics]\label{prop:doob}
The operator $\cP_{\rho,\theta}$ is positive, preserves constants, and has a
spectral gap on a Holder space.  Its stationary probability is
\[
 \nu_{\rho,\theta}=h_{\rho,\theta}\ell_{\rho,\theta}.
\]
Pushed through the same coding map $\chi_\rho$, it is an invariant
preparation of the same deterministic specular billiard $B_\rho$.  Only the
initial path ensemble changes; the collision law and geometry do not.
\end{proposition}
\begin{proof}
The eigenfunction identity gives $\cP_{\rho,\theta}1=1$.  Conjugating the
Ruelle spectral decomposition by $h_{\rho,\theta}$ gives the spectral gap and
stationarity.  The push-forward is invariant because
$B_\rho\circ\chi_\rho=\chi_\rho\circ\sigma$.
\end{proof}

\section{Doob-normalized excess pressure}\label{sec:excess}
Let $F$ be a bounded functional of an $n$-collision future path, measured in
the same impact-action units as the macro preparation.  For $\theta\ne0$ put
\begin{equation}\label{eq:theta-eval}
 \cE_{\rho,r}^{(n)}[F](x)
 =\frac1\theta\log
 \cP_{\rho,\theta}^{n}(e^{\theta F})(x),
 \qquad \theta=\theta_\rho(r).
\end{equation}
At $\theta=0$ use the continuous extension.

\begin{theorem}[Exact excess-pressure identity]\label{thm:excess}
For $\theta\ne0$,
\[
 \cE_{\rho,r}^{(n)}[F](x)
 =\frac1\theta\log
 \frac{
 \cL_{\rho,\theta}^{n}
 (h_{\rho,\theta}e^{\theta F})(x)}
 {e^{n\Lambda_\rho(\theta)}h_{\rho,\theta}(x)}.
\]
The conditional operators are monotone, cash additive, constant preserving,
and satisfy the exact nonlinear tower law under the canonical tilted
preparation.  By \cref{thm:gibbs-conditioning}, this canonical preparation is
the central-window local limit of the finite microcanonical impact
conditioning $K_n^{\leftrightarrow}/(2n)\approx r$.  No finite-$n$
microcanonical Markov property is asserted.
\end{theorem}
\begin{proof}
Iterating \eqref{eq:markov-op} gives the ratio identity.  Positivity and the
linear Markov tower for $\cP_{\rho,\theta}$ imply the nonlinear properties
after exponentiation and logarithm.  The conditioning theorem identifies the
canonical state selected by the macro constraint.
\end{proof}

\section{Geometric response of the ensemble}\label{sec:response}
\begin{theorem}[Joint geometry--preparation regularity]\label{thm:regularity}
For every prescribed finite $N$, the maps
\[
 (\rho,\theta)\mapsto
 h_\rho,\ \Lambda_\rho(\theta),\ h_{\rho,\theta},\
 \ell_{\rho,\theta},\ \nu_{\rho,\theta}
\]
are $C^N$ on compact subsets of $I_\rho\times\R$, provided the obstacle family
is taken $C^{N+2}$.  On compact subsets of
$I_\rho\times(0,1/2)$, $\theta_\rho(r)$ is $C^N$ and
\[
 \partial_r\theta_\rho(r)
 =\frac1{\Lambda_\rho''(\theta_\rho(r))}.
\]
\end{theorem}
\begin{proof}
Use \cref{thm:coding} and operator-norm perturbation of the Ruelle family.
The pressure root $h_\rho$ follows from the implicit-function theorem because
$\partial_hP(-h\widehat\tau_\rho)=-\int\widehat\tau_\rho d\nu<0$.
The inverse function theorem applies to $\Lambda_\rho'$ by
\cref{thm:pressure-duality}.
\end{proof}

\section{Main theorem and scope}\label{sec:main}
\begin{theorem}[Specular first-principles path-ensemble theorem]\label{thm:main}
On the single real specular billiard family $B_\rho$:
\begin{enumerate}[label=(\roman*)]
\item the geometry, collision law, symbolic dynamics, and physical roof are
actual billiard objects, and the marked-period test makes the deformation
geometrically nontrivial;
\item the base path law is selected by the geometry through the
Bowen--Margulis variational principle;
\item a mechanical impact-frequency preparation $r$ determines the unique
conjugate field $\theta_\rho(r)=I_\rho'(r)$;
\item microcanonical finite windows converge quantitatively to the tilted
same-billiard equilibrium state;
\item the Doob-normalized Ruelle operator produces an exact dynamically
consistent excess-pressure theta-evaluation.
\end{enumerate}
No arbitrary Markov matrix, microscopic noise, surrogate collision map, or
independent risk parameter is used in this theorem.
\end{theorem}
\begin{proof}
Combine \cref{prop:no-eclipse,thm:coding,thm:marked-period,thm:BM,thm:pressure-duality,thm:gibbs-conditioning,prop:doob,thm:excess}.
\end{proof}

\section*{Scope and review status}
The theorem concerns the non-wandering set of a no-eclipse open billiard and
its geometry-selected equilibrium preparations.  It is not a Liouville/SRB
theorem for a compact periodic Sinai table.  Independent line-by-line review
has not yet occurred.

\nocite{*}
\bibliographystyle{plain}
\bibliography{references}
\end{document}
